%==============================================================================
% 1 INTRODUCTION
% Plano: reality gap nas plataformas de futebol de robos e, em geral, nos gemeos de
% robos moveis; a camada de planta de um gemeo (mapa comando -> movimento); dados
% operacionais que a equipe ja grava; a pergunta (o que esses dados identificam e o
% que nao identificam); seis contribuicoes; o artigo do SBR como versao preliminar.
% Fonte: SBR sec. I (reescrever, nao copiar); plano §2 (tese); refs de references.bib.
% Contribuicoes candidatas (fechar com a orientadora):
%  (1) identificabilidade como propriedade do regime e da realimentacao (C1, C1b, r3);
%  (2) mecanismo do vies do yaw: acoplamento translacao -> guinada (C1b + MIMO/WP2);
%  (3) regra de triagem antes do ajuste (excitation_per_dof, analyze_omega_excitation);
%  (4) protocolo de fidelidade observacional pre-registrado, agnostico ao simulador,
%      com condicao inicial justa (harness, C2b, C2b2);
%  (5) comparacao com fisica e simuladores da liga sob o mesmo protocolo (C2a, C2b, C2b2);
%  (6) plausibilidade em malha fechada (C4, B3).
% Obrigatorio: "A preliminary version of this work appeared in \cite{sbr2026preliminary};
% this article adds ...".
%==============================================================================
\section{Introduction}
\label{sec:intro}

\pending{motivation, question, contributions, relation to the preliminary version \cite{sbr2026preliminary}.}
